\documentclass[UTF8,fontset=none,11pt,a4paper]{ctexart}
\usepackage{ipara-handbook}
\usepackage{graphicx}
\handbooksetup{NET-03 可靠传输}{408 · 计算机网络 NET-03}{Seq · ACK · Timer · Window}
\providecommand{\tightlist}{\setlength{\itemsep}{0pt}\setlength{\parskip}{0pt}}
\providecommand{\passthrough}[1]{#1}
\begin{document}
\handbooktitle{可靠传输}{用有限状态驯服丢失、损坏与乱序}{408 · NET-03}
\section{本册定位}\label{0-ux672cux518cux5b9aux4f4d}

本 Topic
回答：下层只提供可能损坏、丢失、重复和乱序的交付时，通信端点怎样共同制造``正确且不重复''的可靠服务？

本册拥有 sequence number、ACK、timer、retransmission、sender/receiver
window、Stop-and-Wait、GBN、SR，以及窗口利用率与序号空间约束。

本册建立通用可靠传输机制，不把它等同于 TCP。TCP 如何把 byte
stream、连接状态与这些机制组合，见\href{../06_传输层_端点_UDP与TCP状态机/README.md}{传输层与
TCP}；\passthrough{\lstinline!cwnd!}
不属于本册，见\href{../07_拥塞_共享资源与反馈控制/README.md}{拥塞控制}。

\section{独立阅读词典：可靠到底保证什么}

这里的“可靠”不是承诺链路永不出错，而是让接收应用看到的交付满足题目指定的
顺序、无重复和完整性要求。下表固定本册所有流程图的词义：

\begin{handbooktable}{L{2.2cm} L{4.4cm} Y}
\toprule
术语 & 本册中的可操作定义 & 不要混成什么 \\
\midrule
sequence number / ACK & 序号给数据定位；ACK 反馈“某序号已收到”或“此前连续前缀已收到” & ACK 不是对方一定收到的全局证明 \\
cumulative ACK & 一个值确认此前连续数据；individual ACK 只确认指定数据 & 不能把两种 ACK 记法混用 \\
timer / retransmission & 等待确认的本地时钟；到期后再次发送同一数据 & 超时不等于已证明原数据丢失 \\
window & 当前允许发送或接受的序号范围，决定能有多少数据在途 & 不是文件的固定分块大小 \\
GBN / SR & GBN 失序少缓存、缺口后批量回退；SR 缓存失序、只重传缺失帧 & 不是两种不同的可靠目标 \\
duplicate suppression & 用序号和接收状态保证重传副本不再次向上交付 & 不是阻止网络中出现副本 \\
reliable delivery & 让应用看到题设要求的顺序、无重复和完整性 & 不是承诺链路永不出错 \\
\bottomrule
\end{handbooktable}

因此，“超时”只能推出一个可执行动作：重传或进入题设指定的失败分支；它不能单独
证明原数据丢失，因为 ACK 也可能丢失，或报文仍在传播。

\section{根本问题：沉默不能说明发生了什么}\label{1-ux6839ux672cux95eeux9898ux6c89ux9ed8ux4e0dux80fdux8bf4ux660eux53d1ux751fux4e86ux4ec0ux4e48}

发送一个数据单元后没有收到反馈，至少存在三种解释：

\begin{enumerate}
\def\labelenumi{\arabic{enumi}.}
\tightlist
\item
  数据丢失；
\item
  数据到达，但 ACK 丢失；
\item
  数据或 ACK 仍在路上。
\end{enumerate}

端点没有全局观察者，不能直接知道网络内部发生了什么。可靠传输只能通过有限的报文和本地时间推断状态：

\[
\boxed{
\text{Error detection}
\to \text{Identity}
\to \text{Feedback}
\to \text{Timeout}
\to \text{Retransmission}
\to \text{Duplicate suppression}
}
\]

可靠不是某一个字段的性质，而是发送方与接收方状态机形成的闭环。

\section{五个最小机制为什么缺一不可}\label{2-ux4e94ux4e2aux6700ux5c0fux673aux5236ux4e3aux4ec0ux4e48ux7f3aux4e00ux4e0dux53ef}

\subsection{Error
detection：先知道``这个副本不能信''}\label{21-error-detectionux5148ux77e5ux9053ux8fd9ux4e2aux526fux672cux4e0dux80fdux4fe1}

校验和或 CRC 让接收方检测部分传输错误。检测失败通常触发丢弃或
NAK，但它不能恢复原数据。

\subsection{Sequence
number：给传输实例一个可比较身份}\label{22-sequence-numberux7ed9ux4f20ux8f93ux5b9eux4f8bux4e00ux4e2aux53efux6bd4ux8f83ux8eabux4efd}

重传会产生多个内容相同的副本。若没有序号，接收方无法判断``这是新数据还是旧数据重传''。序号还为顺序和累计确认提供坐标系。

\subsection{ACK：把接收端知识反馈给发送端}\label{23-ackux628aux63a5ux6536ux7aefux77e5ux8bc6ux53cdux9988ux7ed9ux53d1ux9001ux7aef}

ACK
的语义必须先写清：确认单个序号、累计确认到某个序号，还是声明下一个期望序号。不同教材或协议可能采用不同记法，不能只凭
\passthrough{\lstinline!ACK(n)!} 猜。

\subsection{Timer：把无限等待变成可执行分支}\label{24-timerux628aux65e0ux9650ux7b49ux5f85ux53d8ux6210ux53efux6267ux884cux5206ux652f}

发送方无法等到``确定丢失''，只能在等待超过阈值后推定本次尝试失败。Timer
太短会产生伪重传，太长会延迟恢复。

\subsection{Retransmission + duplicate
suppression}\label{25-retransmission--duplicate-suppression}

超时重传解决可能丢失；序号与接收状态消除重传造成的重复交付。二者共同把``至少一次尝试''收缩为``至多一次向上交付''。

\section{Stop-and-Wait：最小闭环}\label{3-stop-and-waitux6700ux5c0fux95edux73af}

发送方状态：

\begin{lstlisting}
Ready(n)
-> send packet(n), start timer
-> WaitACK(n)
   -> valid ACK(n): stop timer, n toggles, Ready(next)
   -> timeout: resend packet(n), restart timer
\end{lstlisting}

接收方状态：

\begin{lstlisting}
Expect(n)
-> valid packet(n): deliver once, send ACK(n), Expect(next)
-> duplicate/out-of-order: do not deliver, repeat appropriate ACK
-> corrupted: discard or send feedback according to protocol
\end{lstlisting}

1-bit 序号足以支持 Stop-and-Wait，因为任意时刻只需区分当前新 packet
与上一个 packet 的迟到/重传副本。

\subsection{正确性不变量}\label{31-ux6b63ux786eux6027ux4e0dux53d8ux91cf}

\begin{itemize}
\tightlist
\item
  发送方收到足以推进状态的确认前，不复用当前序号；
\item
  接收方只把当前期望序号向上交付一次；
\item
  ACK 丢失可以造成重复传输，但不能造成重复交付。
\end{itemize}

\subsection{性能失败}\label{32-ux6027ux80fdux5931ux8d25}

若数据发送时间为 \(T_D\)、单向传播时延为 \(T_P\)，并把 \(RTT\) 简化为
\(2T_P\)，忽略 ACK 发送与处理：

\[
U\approx\frac{T_D}{T_D+RTT}.
\]

当
\(RTT\gg T_D\)，链路大部分时间空闲。可靠性已经成立，但吞吐能力没有被利用，于是需要流水线。

\section{Sliding
Window：把等待时间变成在途数据}\label{4-sliding-windowux628aux7b49ux5f85ux65f6ux95f4ux53d8ux6210ux5728ux9014ux6570ux636e}

窗口不是``批量发送''的口诀，而是对允许未确认序号集合的约束：

\[
\boxed{
\text{Past ACKed}
\mid \text{Sent unACKed}
\mid \text{Usable unsent}
\mid \text{Not yet allowed}
}
\]

ACK 到达使左边界右移，释放新的序号。窗口同时承担：

\begin{itemize}
\tightlist
\item
  correctness：限制哪些序号当前合法；
\item
  flow of work：允许多少数据在反馈返回前继续发送；
\item
  bounded state：限定发送方、接收方和缓冲区必须维护多少状态。
\end{itemize}

若窗口含 \(W\) 个等长 frame，ACK 发送时间为
\(T_{ACK}\)，理想无差错利用率近似：

\[
U=\min\left(1,\frac{W T_D}{T_D+RTT+T_{ACK}}\right).
\]

公式必须从时间线生成。题目忽略 ACK 发送时间时才去掉 \(T_{ACK}\)。

\section{GBN：用接收端简单换取批量回退}\label{5-gbnux7528ux63a5ux6536ux7aefux7b80ux5355ux6362ux53d6ux6279ux91cfux56deux9000}

GBN 的接收窗口通常为 1：只接收当前期望 frame，失序 frame
不缓存。发送方允许多个未确认 frame 在途，并使用累计确认。

\subsection{生命周期}\label{51-ux751fux547dux5468ux671f}

\begin{lstlisting}
sender sends base ... nextseq-1
-> receiver accepts only expected frame
-> cumulative ACK advances base
-> if oldest unACKed timer expires
-> resend base and every later outstanding frame
\end{lstlisting}

\subsection{为什么只需一个主计时器}\label{52-ux4e3aux4ec0ux4e48ux53eaux9700ux4e00ux4e2aux4e3bux8ba1ux65f6ux5668}

累计确认使最老未确认 frame 成为窗口能否前进的阻塞点。其超时后，后续在途
frame 即使曾正确到达也可能因接收方不缓存而需要一起重传。

\subsection{序号空间约束}\label{53-ux5e8fux53f7ux7a7aux95f4ux7ea6ux675f}

若序号字段为 \(k\) bit，在 408 常用模型中：

\[
W_T\le 2^k-1,
\qquad W_R=1.
\]

必须保留至少一个序号不在当前发送窗口中，才能让新一轮窗口与旧副本保持可区分。

\section{SR：用更多状态换取精准恢复}\label{6-srux7528ux66f4ux591aux72b6ux6001ux6362ux53d6ux7cbeux51c6ux6062ux590d}

SR 的接收方可以缓存窗口内失序 frame，并逐个确认；发送方通常为每个未确认
frame 维护独立计时状态。丢一个 frame 时只重传所缺 frame。

\subsection{生命周期}\label{61-ux751fux547dux5468ux671f}

\begin{lstlisting}
frame n arrives within receive window
-> if first valid copy: buffer and ACK n
-> if n equals receive base: deliver n and consecutive buffered frames
-> slide receive window over delivered prefix
\end{lstlisting}

接收方可以``已收到但暂不能交付''。这说明 reception state 与 delivery
state 不是同一件事。

\subsection{序号复用为什么会产生歧义}\label{62-ux5e8fux53f7ux590dux7528ux4e3aux4ec0ux4e48ux4f1aux4ea7ux751fux6b67ux4e49}

序号是有限循环空间。若发送窗口和接收窗口覆盖范围过大，一个迟到的旧 frame
可能落入新一轮接收窗口，并被误认为新 frame。

通用不歧义条件是：

\[
W_T+W_R\le 2^k.
\]

常用对称窗口下：

\[
W_T=W_R\le 2^{k-1}.
\]

这个约束不是为了``方便计算''，而是为了让接收方仅凭有限序号和当前窗口就能判定身份。

\begin{center}
\resizebox{0.96\linewidth}{!}{\providecolor{themebg}{HTML}{FAFAF7}
\providecolor{themefg}{HTML}{111111}
\providecolor{themegray}{HTML}{666666}
\providecolor{themecurve}{HTML}{1D63B8}
\providecolor{themealert}{HTML}{C53030}
\providecolor{themeamber}{HTML}{B86E00}
\begin{tikzpicture}[font=\scriptsize,color=themefg,text=themefg,>=Stealth]
\tikzset{
  seq/.style={draw=themegray,minimum width=7.5mm,minimum height=7.5mm,inner sep=0pt},
  send/.style={seq,draw=themecurve,fill=themecurve!10!themebg},
  recv/.style={seq,draw=themeamber,fill=themeamber!10!themebg},
  overlap/.style={seq,draw=themealert,fill=themealert!12!themebg},
  note/.style={rounded corners=3pt,draw=themegray,fill=themefg!3!themebg,inner xsep=6pt,inner ysep=5pt,align=left}
}
\node[anchor=west,font=\bfseries\large] at (0,9.2) {Stop-and-Wait、GBN、SR：窗口上限来自“旧副本与新一轮序号不能混淆”};
\node[anchor=west,text=themegray] at (0,8.68) {例用 3-bit 序号空间 $0\ldots7$；先锁 ACK、缓存和 Timer 语义，再把窗口解释为当前合法身份集合。};

% SAW
\node[anchor=west,font=\bfseries,text=themecurve] at (0,7.65) {A. Stop-and-Wait};
\foreach \i in {0,...,7}{\pgfmathsetmacro{\xx}{2.2+0.78*\i}\ifnum\i=0\node[send] at (\xx,7.05) {0};\else\node[seq] at (\xx,7.05) {\i};\fi}
\node[note,anchor=west,text width=6.8cm] at (9.0,7.05) {任意时刻只允许 1 个未确认数据实例；1 bit 就足以区分“当前新帧”和“上一个序号的迟到/重传副本”。};

% GBN
\node[anchor=west,font=\bfseries,text=themecurve] at (0,5.85) {B. GBN：$W_R=1$，发送窗口至少留 1 个序号在外};
\foreach \i in {0,...,7}{\pgfmathsetmacro{\xx}{2.2+0.78*\i}\ifnum\i<7\node[send] at (\xx,5.25) {\i};\else\node[seq,draw=themealert] at (\xx,5.25) {7};\fi}
\node[anchor=north,text=themecurve] at (4.54,4.73) {$W_T=7=2^3-1$ 可取满};
\node[anchor=north,text=themealert] at (7.66,4.73) {至少一个序号不在本轮发送窗口};
\node[note,anchor=west,text width=6.8cm] at (9.0,5.25) {接收端只接受下一期望序号，失序帧通常丢弃；累计 ACK 压缩连续前缀。若 8 个序号全进发送窗口，回绕后新 0 与仍可能存在的旧 0 无法靠序号区分。};

% SR
\node[anchor=west,font=\bfseries,text=themecurve] at (0,3.55) {C. SR：发送/接收窗口都大，必须避免新旧接收区间重叠};
\foreach \i in {0,...,7}{\pgfmathsetmacro{\xx}{2.2+0.78*\i}\ifnum\i<4\node[recv] (old\i) at (\xx,2.95) {\i};\else\node[seq] at (\xx,2.95) {\i};\fi}
\node[anchor=north,text=themeamber] at (3.37,2.43) {旧接收窗口 $[0,3]$};
\foreach \i in {0,...,7}{\pgfmathsetmacro{\xx}{2.2+0.78*\i}\ifnum\i>3\node[send] at (\xx,1.75) {\i};\else\node[seq] at (\xx,1.75) {\i};\fi}
\node[anchor=north,text=themecurve] at (6.49,1.23) {安全的新窗口 $[4,7]$：与旧窗口不重叠};
\node[note,anchor=west,text width=6.8cm] at (9.0,2.35) {对称 SR 取 $W_T=W_R=4=2^{3-1}$ 时，序号空间恰好分成两半。旧帧 0 即使迟到，也不属于新合法窗口 $[4,7]$。};

% unsafe overlap minimal counterexample
\node[anchor=west,font=\bfseries,text=themealert] at (0,.45) {若 SR 窗口扩大到 5：};
\foreach \i in {0,...,7}{\pgfmathsetmacro{\xx}{2.2+0.78*\i}\ifnum\i<5\node[recv] at (\xx,-.2) {\i};\else\node[seq] at (\xx,-.2) {\i};\fi}
\node[anchor=west] at (6.8,-.2) {$\to$ 滑动 5 后新窗口为 $[5,6,7,0,1]$};
\node[overlap] at (11.25,-.2) {0}; \node[overlap] at (12.05,-.2) {1};
\node[anchor=west,text=themealert] at (12.7,-.2) {旧 0/1 可落入新窗口，身份歧义};
\node[anchor=west,font=\bfseries,text=themecurve] at (0,-1.55) {$\boxed{\text{GBN: }W_T\le2^k-1\qquad\text{SR: }W_T+W_R\le2^k\quad(\text{对称时 }W\le2^{k-1})}$};
\node[anchor=west,font=\scriptsize,text=themegray] at (0,-2.15) {窗口的性能需求由 RTT/发送时间产生；这里画的是协议身份安全上限，两类约束不能互相替代。};
\end{tikzpicture}
}
\end{center}
图用同一个 3-bit 序号空间对照 Stop-and-Wait、GBN 与 SR，并直接构造 SR 窗口过大时旧帧重新落入新接收窗口的身份歧义。

\section{GBN 与 SR
的真正权衡}\label{7-gbn-ux4e0e-sr-ux7684ux771fux6b63ux6743ux8861}

\begin{longtable}[]{@{}lll@{}}
\toprule\noalign{}
维度 & GBN & SR \\
\midrule\noalign{}
\endhead
\bottomrule\noalign{}
\endlastfoot
接收失序 frame & 丢弃 & 缓存 \\
ACK & 通常累计 & 通常逐 frame \\
超时恢复 & 从缺口开始批量重传 & 只重传具体缺失 frame \\
接收端状态 & 小 & 大 \\
发送端 timer & 通常最老 frame 一个主 timer & 每个未确认 frame
独立状态 \\
高误码链路代价 & 重传浪费较大 & 状态和实现复杂度较大 \\
\end{longtable}

SR 不是无条件更好。链路很可靠、窗口较小或接收资源受限时，GBN
的简单性可能更值钱。

\section{408 状态机覆盖：发送端与接收端双账本}

任何 ARQ 流程都用同一张事件表推进：

\begin{handbooktable}{L{2.4cm} L{2.8cm} Y Y}
\toprule
事件 & 报文/字段 & 发送端状态变化 & 接收端状态变化\\
\midrule
发送新 frame & seq、payload、checksum & 占用窗口槽；启动相应 timer & 无\\
正确到达 & seq 在/不在接收窗口 & 等待 ACK & 接收/缓存/交付；产生 ACK\\
损坏或失序 & checksum/seq 不满足 & 等待 ACK 或超时 & 丢弃、重复 ACK 或 NAK（按题设）\\
ACK 到达 & ACK 语义必须先声明 & 标记确认；连续前缀满足才滑窗 & 无\\
timer 到期 & 对应 seq 或 oldest unACKed & Stop-and-Wait/GBN/SR 按各自策略重传 & 无\\
\bottomrule
\end{handbooktable}

\subsection{Stop-and-Wait 的四个异常分支}

必须分别检查 data loss、data corruption、ACK loss、ACK corruption/delay。发送方超时后都可能重传；接收方用期望序号保证重复副本只回 ACK、不重复 deliver。停止条件不是“某个副本到达”，而是 sender 获得足以安全复用序号的反馈。

\subsection{GBN：base、nextseqnum 与唯一主 timer}

发送端允许
\[
base\le seq<base+W_T.
\]
ACK 的累计语义只推进连续确认前缀。收到旧 ACK 不回退 base；最老未确认 frame 超时则重传区间 $[base,nextseqnum)$。接收端窗口为 1，失序 frame 不缓存，并重复报告当前连续前缀。每个事件后检查
\[
base\le nextseqnum\le base+W_T.
\]

\subsection{SR：接收、缓存、交付是三种状态}

发送端逐 frame 维护 ACK 与 timer；接收端对窗口内首个有效副本缓存并选择确认。只有 receive base 已到达时，才把它和随后连续缓存项一起 deliver 并滑动窗口。迟到但已落在旧窗口的副本可重复确认却不得再交付；窗口外未来序号按题设丢弃。序号空间必须满足 $W_T+W_R\le 2^k$，否则仅凭 seq 无法区分旧副本与新循环。

\begin{warnbox}{停止类比}
本册的 GBN/SR 是通用 ARQ 模型。TCP 可以调用累计 ACK、失序缓存和选择确认等机制，但不能被整体贴成 GBN 或 SR；TCP 的 byte sequence、connection 和 flow control 由 NET06 Own。
\end{warnbox}

\section{计算与反例生成：从反馈周期到序号空间}

\subsection{一个完整周期不能漏掉 ACK 与处理阶段}

设数据帧发送时延为 $T_f$、ACK 发送时延为 $T_a$、单向传播时延为 $\tau$、接收处理时延为 $T_p$。停止等待从开始发送一帧到可以安全发送下一帧的周期为
\[
T_{cycle}=T_f+\tau+T_p+T_a+\tau
=T_f+2\tau+T_p+T_a,
\qquad
U_{SW}=\frac{T_f}{T_{cycle}}.
\]
只有题设明确忽略 ACK 与处理开销时，才化为
\[
U_{SW}=\frac{1}{1+2a},\qquad a=\frac{\tau}{T_f}.
\]
这说明提高链路速率会缩短 $T_f$，却不缩短 $\tau$；若窗口仍为 1，$a$ 反而变大，利用率可能下降。超时阈值应覆盖正常反馈周期并留有余量：过短产生伪重传，过长延迟真实丢失恢复。链路层 RTT 较稳定时可用固定余量，端到端 TCP 的动态 RTT 估计由 NET06 拥有。

\subsection{流水线满载阈值与序号位数反推}

发送窗口为 $W$ 时，理想无差错利用率为
\[
U(W)=\min\!\left(1,\frac{W T_f}{T_f+2\tau+T_p+T_a}\right).
\]
因此满载所需的最小整数窗口为
\[
W_{\mathrm{sat}}
=\left\lceil\frac{T_f+2\tau+T_p+T_a}{T_f}\right\rceil.
\]
忽略 $T_p,T_a$ 时，$W_{\mathrm{sat}}=\lceil1+2a\rceil$。再把物理阈值接到协议上限：
\[
\begin{aligned}
\text{GBN:}\quad &2^n-1\ge W_{\mathrm{sat}}
&&\Longrightarrow&& n\ge\left\lceil\log_2(W_{\mathrm{sat}}+1)\right\rceil,\\
\text{SR, }W_T=W_R:\quad &2^{n-1}\ge W_{\mathrm{sat}}
&&\Longrightarrow&& n\ge1+\left\lceil\log_2W_{\mathrm{sat}}\right\rceil.
\end{aligned}
\]
顺序必须是“先由链路时间线求窗口需求，再由协议不歧义约束求序号位数”，不能直接把某个 $2^n$ 公式当利用率。

\subsection{ACK 丢失为什么在 GBN 与 SR 中代价不同}

GBN 的 ACK 是连续前缀摘要。若 ACK 3 丢失而 ACK 5 到达，ACK 5 已包含“0--5 都按序收到”的知识，发送方可以直接推进到 6；较大的累计 ACK 覆盖较小 ACK。SR 的 ACK 通常是逐帧状态：ACK 3 丢失不能由 ACK 4 或 ACK 5 自动证明 frame 3 已收到，因此 frame 3 的 timer 仍可能到期并产生一次冗余重传。接收方收到已交付窗口之前的旧 frame 时，必须重新发送该 frame 的 ACK，否则发送方可能永远无法关闭对应 timer。

这不是“SR 的 ACK 更差”，而是反馈信息结构不同：
\[
\boxed{\text{GBN：压缩连续前缀，恢复粒度粗}}
\qquad
\boxed{\text{SR：表达任意接收模式，状态与定时器更多}}.
\]

\subsection{发送窗口与接收窗口不要求同步}

发送窗口在最早未确认项获得足够确认时滑动；接收窗口在其左端起的连续帧已经收到并可交付时滑动。两个条件分别发生在 ACK 回程与 data 到达事件上，所以窗口位置天然可以不同步。SR 的发送窗口内部还可能出现“前项未确认、后项已确认”的洞；窗口左边界不能越过洞，但相应的后项无需重传。

接收 frame 时要分三类：
\begin{itemize}
\tightlist
\item 在当前接收窗口内：首次有效副本则缓存并 ACK；只有补齐左端连续前缀时才按序交付并滑窗；
\item 在接收窗口之前：数据已处理，丢弃副本但重复 ACK；
\item 在接收窗口之后：尚不合法，按题设丢弃。
\end{itemize}

\subsection{有限序号的隐藏前提：旧副本寿命必须有界}

Stop-and-Wait 用 1 bit 序号，只能区分“当前期望的新帧”和“上一帧的重传”。它隐含旧副本不会跨越多个完整序号循环后仍从信道中出现。短而有界的链路模型通常满足这一前提；若网络允许报文任意长时间滞留，仅靠 0/1 交替不能防止极迟旧副本被当成新数据。窗口不重叠约束解决一个反馈周期内的新旧歧义，报文最大寿命与更大序号空间则解决跨周期身份复用风险。

\section{来源覆盖附录：NET-03 逐 Source 核销}

\begingroup\small
\begin{handbooktable}{L{4.2cm} L{3.0cm} Y}
\toprule
Source & 本册承接节点 & 核销结论\\
\midrule
\texttt{datalink/stop-and-wait} & 最小闭环/反馈周期 & 四异常分支、1-bit 序号、timeout 权衡、含 ACK/处理时延的利用率与满载反推已展开\\
\texttt{datalink/gbn} & 累积确认/批量回退 & $W_R=1$、$W_T\le2^n-1$、唯一主 timer、乱序丢弃、ACK 覆盖与重传区间已展开\\
\texttt{datalink/sr} & 逐帧确认/选择重传 & 缓存与按序交付、逐 frame timer、三类接收位置、$W_T+W_R\le2^n$ 与旧帧反例已展开\\
\texttt{datalink/sliding-window-compare} & 统一窗口模型 & 四区域、两端异步滑动、分段利用率、三协议上限/代价/适用条件及传输层边界已展开\\
\bottomrule
\end{handbooktable}
\endgroup

\section{概念边界}\label{8-ux6982ux5ff5ux8fb9ux754c}

\begin{longtable}[]{@{}lclll@{}}
\toprule\noalign{}
概念 A & ≠ & 概念 B & 真正区别与题目信号 & 混淆后果 \\
\midrule\noalign{}
\endhead
\bottomrule\noalign{}
\endlastfoot
Detection & ≠ & Recovery & 检错指出异常；反馈/超时/重传恢复 & 认为 CRC
自带重传 \\
Receive & ≠ & Deliver & frame 可到达并缓存，但未必按序交给上层 & SR
状态推演错误 \\
ACK loss & ≠ & Data loss & 前者接收方已有数据；后者没有 &
漏掉重复抑制分支 \\
Window size & ≠ & Sequence space &
窗口是当前合法集合；序号空间是循环身份全集 & GBN/SR 最大窗口算错 \\
Reliability & ≠ & Flow control & 前者保护数据交付；后者保护接收缓冲 & 把
\passthrough{\lstinline!rwnd!} 当作可靠性条件 \\
Flow control & ≠ & Congestion control & 前者看 receiver；后者看 network
path & 把 \passthrough{\lstinline!rwnd!}、\passthrough{\lstinline!cwnd!}
合并成一个含义 \\
Timeout & ≠ & Proof of loss & 超时是基于本地时钟的推断，不是全局事实 &
无法解释伪重传 \\
\end{longtable}

\section{Verification：状态题怎样提前发现错误}\label{9-verificationux72b6ux6001ux9898ux600eux6837ux63d0ux524dux53d1ux73b0ux9519ux8bef}

每一步检查四个集合：

\begin{enumerate}
\def\labelenumi{\arabic{enumi}.}
\tightlist
\item
  已确认/已交付的连续前缀；
\item
  已发送未确认集合；
\item
  已接收但未交付缓存；
\item
  当前窗口允许的新序号。
\end{enumerate}

再检查三个不变量：

\begin{itemize}
\tightlist
\item
  同一数据不能向上交付两次；
\item
  窗口只跨过连续满足条件的前缀；
\item
  当前接受区间不能与仍可能出现的旧副本产生身份歧义。
\end{itemize}

\section{做题调用协议}\label{10-ux505aux9898ux8c03ux7528ux534fux8bae}

\begin{enumerate}
\def\labelenumi{\arabic{enumi}.}
\tightlist
\item
  先写 ACK 语义和接收方是否缓存失序 frame；
\item
  画 sender/receiver 两条时间线，标出发送、到达、ACK、超时；
\item
  每个事件后更新 base、nextseq、buffer、timer；
\item
  窗口题先判断 GBN 还是 SR，再写对应不歧义条件；
\item
  利用率题从一个 frame 的反馈周期推导，不先套公式；
\item
  有误码/丢包时同时算重传数据量和状态成本；
\item
  最后用``是否可能重复交付''攻击结果。
\end{enumerate}

\section{贯穿母例：ACK
丢失为什么不会破坏可靠性}\label{11-ux8d2fux7a7fux6bcdux4f8back-ux4e22ux5931ux4e3aux4ec0ux4e48ux4e0dux4f1aux7834ux574fux53efux9760ux6027}

Stop-and-Wait 中，发送方发送 frame 0；接收方正确接收、交付并回 ACK 0，但
ACK 丢失：

\begin{lstlisting}
sender remains WaitACK(0)
-> timer expires
-> sender retransmits frame 0
-> receiver is already Expect(1)
-> recognizes duplicate 0, does not deliver again
-> repeats ACK 0
-> sender advances to frame 1
\end{lstlisting}

重传保证 progress，序号与接收状态保证 safety。可靠性的两面分别是：

\[
\boxed{\text{eventually make progress}}
\qquad
\boxed{\text{never deliver the wrong instance twice}}.
\]

\section{高频 First
Divergence}\label{12-ux9ad8ux9891-first-divergence}

\begin{itemize}
\tightlist
\item
  看见丢包直接重传：没有说明谁、凭什么知道丢失；
\item
  ACK 到达就把窗口跳到该序号之后：没有确认 ACK 是累计还是选择确认；
\item
  SR 收到失序 frame 立即交付：混淆 buffer 与 in-order delivery；
\item
  GBN 只重传超时的一个 frame：把恢复策略套成 SR；
\item
  只背 \(2^k-1\) 或 \(2^{k-1}\)：没有从旧/新副本不歧义生成约束；
\item
  利用率超过 1：没有执行饱和截断或时间线有重复计数。
\end{itemize}

\section{一页压缩与复原问题}\label{13-ux4e00ux9875ux538bux7f29ux4e0eux590dux539fux95eeux9898}

\[
\boxed{
\text{Detect}
\to \text{Number}
\to \text{Acknowledge}
\to \text{Time out}
\to \text{Retransmit}
\to \text{Suppress duplicate}
}
\]

\begin{enumerate}
\def\labelenumi{\arabic{enumi}.}
\tightlist
\item
  为什么``没收到 ACK''不能等价为``数据丢了''？
\item
  Stop-and-Wait 为什么 1-bit 序号足够？
\item
  窗口怎样把传播等待转化为在途数据？
\item
  GBN 和 SR 各把复杂度放在哪一端？
\item
  SR 的序号空间约束怎样从迟到旧 frame 推出？
\end{enumerate}

\section{来源与校正说明}\label{14-ux6765ux6e90ux4e0eux6821ux6b63ux8bf4ux660e}

\begin{itemize}
\tightlist
\item
  归档笔记《链路层-流量控制与可靠传输》《公式汇总》提供 GBN/SR
  考点和旧计算模型；
\item
  本册把可靠传输提升为跨层可复用机制，避免被旧目录永久绑定在链路层；
\item
  ACK
  记号在不同教材中可能不同，正文要求先声明语义，不把单一记法冒充协议普遍事实。
\end{itemize}
\end{document}
